\documentclass[11pt,a4paper]{article}
\usepackage[utf8]{inputenc}
\usepackage[T1]{fontenc}
\usepackage[brazilian]{babel}
\usepackage{geometry}
\geometry{top=1.7cm, bottom=1.7cm, left=1.8cm, right=1.8cm}
\usepackage{amsmath, amssymb}
\usepackage{xcolor}
\usepackage{titlesec}
\usepackage{enumitem}
\usepackage{booktabs}
\usepackage{tabularx}
\usepackage{tcolorbox}
\usepackage{hyperref}
\usepackage{microtype}
\usepackage{newtxtext,newtxmath}

% Paleta Institucional IFES
\definecolor{ifesgreen}{RGB}{24, 138, 63}
\definecolor{ifesdark}{RGB}{14, 88, 39}
\definecolor{ifesred}{RGB}{200, 30, 30}
\definecolor{darkgray}{RGB}{45, 50, 58}
\definecolor{lightgray}{RGB}{246, 248, 249}
\definecolor{boxbg}{RGB}{240, 247, 242}
\definecolor{bordergreen}{RGB}{175, 215, 190}

\hypersetup{
    colorlinks=true,
    linkcolor=ifesgreen,
    urlcolor=ifesgreen
}

\titleformat{\section}{\color{ifesgreen}\large\bfseries}{\thesection}{0.8em}{}[\color{ifesgreen}\titlerule]
\titleformat{\subsection}{\color{ifesdark}\normalsize\bfseries}{\thesubsection}{0.8em}{}

\tcbset{
    boxrule=0.8pt,
    arc=4pt,
    left=7pt, right=7pt, top=5pt, bottom=5pt
}

\begin{document}

% Cabeçalho
\begin{center}
    {\small\textbf{INSTITUTO FEDERAL DO ESPÍRITO SANTO -- IFES | CAMPUS SERRA}}\\[2pt]
    {\small\textbf{PROGRAMA DE PÓS-GRADUAÇÃO EM ENGENHARIA DE CONTROLE E AUTOMAÇÃO (PROPECAUT)}}\\[4pt]
    \rule{\textwidth}{1.2pt}\\[5pt]
    {\large\textbf{\color{ifesgreen}ESPECIFICAÇÃO MATEMÁTICA DOS MODELOS DO CAPÍTULO 3}}\\[2pt]
    {\small\textbf{Subsistema Elétrico, Armazenamento Eletroquímico e Dinâmica Eletrotérmica de Cargas de IA}}\\[4pt]
    \rule{\textwidth}{0.6pt}
\end{center}

\vspace{-6pt}
\begin{center}
    \small
    \begin{tabularx}{\textwidth}{lX}
        \textbf{Dissertação:} & Controle e Gerenciamento Térmico-Energético Computacional de um Modelo de Data Center Orbital \\
        \textbf{Mestrando / Orientador:} & Dafnhê Thor C. M. Barros \;|\; Prof. Dr. Flávio Barcelos Braz da Silva \\
        \textbf{Nó de Referência:} & Módulo Satelital de $50\text{ kW}$ brutos ($40\text{ kW}$ úteis de IA) em Órbita SSO \textit{Dawn-Dusk} ($700\text{ km}$)
    \end{tabularx}
\end{center}

\vspace{2pt}

% ==============================================================================
% MODELO 1
% ==============================================================================
\begin{tcolorbox}[colback=boxbg, colframe=bordergreen, title=\color{ifesdark}\textbf{MODELO 1: GERAÇÃO SOLAR ORBITAL (ARRANJO FOTOVOLTAICO AM0)}]
\small
\textbf{Finalidade:} Determinar a potência elétrica bruta disponível $P_{pv}(t)$ gerada pelos painéis de tripla junção ($\text{GaInP/GaAs/Ge}$, referência Spectrolab XTJ Prime) sob espectro extraterrestre AM0 ao longo da trajetória orbital.

\vspace{-4pt}
\begin{equation}
    P_{pv}(t) = G(t) \cdot A_{pv} \cdot \eta_{ref} \Big[ 1 - \gamma_T \big( T_{pv}(t) - T_{ref} \big) \Big] \cdot f_{EOL}
    \label{eq:modelo_pv}
\end{equation}
onde a irradiância incidente $G(t)$ segue a função geométrica de eclipse umbral:
\begin{equation}
    G(t) = G_{AM0} \cdot u_{sol}(t) \cdot \cos\theta_{inc}(t), \qquad
    u_{sol}(t) =
    \begin{cases}
        1, & \text{em insolação direta} \\
        0, & \text{em eclipse umbral}
    \end{cases}
    \label{eq:irradiancia}
\end{equation}

\vspace{2pt}
\noindent\textbf{Parâmetros Nominais:} $G_{AM0} = 1361\text{ W/m}^2$; $\eta_{ref} = 0{,}307$ (BOL a $T_{ref} = 28^\circ\text{C}$); $\gamma_T = 0{,}0022\text{ K}^{-1}$; fator de fim de vida $f_{EOL} = 0{,}85$; eclipses máximos $t_{ecl,SSO} = 18{,}8\text{ min}$ (solstício, $\beta_{min}=58{,}3^\circ$) e $t_{ecl,LEO} = 35{,}8\text{ min}$ (cenário comparativo).\\
\textbf{Entradas $\to$ Saída:} Sinal orbital de eclipse $u_{sol}(t)$ e temperatura do painel $T_{pv}(t)$ $\;\longrightarrow\;$ \textbf{Potência gerada $P_{pv}(t)$}.
\end{tcolorbox}

\vspace{4pt}

% ==============================================================================
% MODELO 2
% ==============================================================================
\begin{tcolorbox}[colback=boxbg, colframe=bordergreen, title=\color{ifesdark}\textbf{MODELO 2: BALANÇO DE POTÊNCIA NO BARRAMENTO HVDC (50 kW) E PERDAS GaN}]
\small
\textbf{Finalidade:} Descrever a conservação instantânea de potência no barramento regulado ($400\text{ V} \to 48\text{ V} \to 0{,}8\text{ V}$) e quantificar o calor secundário gerado pelas perdas dos conversores CC-CC em Nitreto de Gálio (GaN).

\vspace{-4pt}
\begin{equation}
    P_{pv}(t) - P_{bat}(t) = \frac{P_{TI}(t)}{\eta_{conv}} + P_{aux}
    \label{eq:balanco_hvdc}
\end{equation}
em que a eficiência em cascata dos estágios GaN e o calor residual da eletrônica de potência são dados por:
\begin{equation}
    \eta_{conv} = \eta_{400\to 48} \cdot \eta_{48\to 0{,}8}, \qquad
    Q_{perdas}(t) = P_{TI}(t) \left( \frac{1}{\eta_{conv}} - 1 \right)
    \label{eq:perdas_gan}
\end{equation}

\vspace{2pt}
\noindent\textbf{Parâmetros Nominais:} Potência nominal do nó $P_{no} = 50\text{ kW}$; carga útil máxima de IA $P_{TI,max} = 40\text{ kW}$; eficiência agregada $\eta_{conv} \in [0{,}95,\; 0{,}97]$; cargas auxiliares $P_{aux} \approx 2{,}0\text{ kW}$ (OBC, ADCS e enlaces ópticos OISL). Convenção: $P_{bat} > 0$ indica carga da bateria e $P_{bat} < 0$ indica descarga.\\
\textbf{Entradas $\to$ Saídas:} $P_{pv}(t)$ e $P_{TI}(t)$ $\;\longrightarrow\;$ \textbf{Fluxo líquido $P_{bat}(t)$} e \textbf{Calor secundário $Q_{perdas}(t)$}.
\end{tcolorbox}

\vspace{4pt}

% ==============================================================================
% MODELO 3
% ==============================================================================
\begin{tcolorbox}[colback=boxbg, colframe=bordergreen, title=\color{ifesdark}\textbf{MODELO 3: DINÂMICA MACRO DO ARMAZENAMENTO ELETROQUÍMICO (BESS)}]
\small
\textbf{Finalidade:} Acompanhar a evolução temporal do Estado de Carga ($SoC$) do banco de baterias de íons de lítio durante eclipses sazonais e transitórios de IA, impondo a restrição energética ao controlador preditivo.

\vspace{-4pt}
\begin{equation}
    \frac{d\,SoC(t)}{dt} =
    \begin{cases}
        \dfrac{\eta_{ch} \, P_{bat}(t)}{E_{nom}}, & \text{se } P_{bat}(t) > 0 \quad (\text{carregamento}) \\[10pt]
        \dfrac{P_{bat}(t)}{\eta_{dis} \, E_{nom}}, & \text{se } P_{bat}(t) < 0 \quad (\text{descarregamento})
    \end{cases}
    \label{eq:soc_dinamica}
\end{equation}
sujeita às restrições operacionais de segurança e longevidade ($DoD_{max} \le 30\%$):
\begin{equation}
    SoC_{min} \le SoC(t) \le SoC_{max}, \qquad \big|P_{bat}(t)\big| \le P_{bat,max}
    \label{eq:soc_restricoes}
\end{equation}

\vspace{2pt}
\noindent\textbf{Parâmetros Nominais:} Capacidade nominal $E_{nom}$ dimensionada para sustentar $t_{ecl,SSO} = 18{,}8\text{ min}$ respeitando $SoC_{min} = 0{,}70$ e $SoC_{max} = 0{,}95$; rendimentos $\eta_{ch} = \eta_{dis} = 0{,}98$ (eficiência de ciclo $\approx 96\%$).\\
\textbf{Entradas $\to$ Saída:} Potência líquida $P_{bat}(t)$ $\;\longrightarrow\;$ \textbf{Estado de Carga $SoC(t)$} (restrição de estado do TEMS).
\end{tcolorbox}

\vspace{4pt}

% ==============================================================================
% MODELO 4
% ==============================================================================
\begin{tcolorbox}[colback=boxbg, colframe=bordergreen, title=\color{ifesdark}\textbf{MODELO 4: DINÂMICA ELETROTÉRMICA DAS CARGAS DE IA (GERAÇÃO DE CALOR E DVFS)}]
\small
\textbf{Finalidade:} Conectar o despacho computacional de IA e a atuação do controle DVFS ($f_{clk}, V_{core}$) à injeção de calor na estrutura do satélite. \textbf{Este modelo é a interface direta de acoplamento com o Capítulo 4.}

\vspace{-4pt}
\begin{equation}
    P_{TI}\big(f_{clk}, V_{core}, T_j\big) = N_{chips} \Big[ \underbrace{\alpha(t) \, C_{eff} \, V_{core}^2 \, f_{clk}}_{\text{Potência Dinâmica } (P_{dyn})} + \underbrace{V_{core} \, I_{leak,0} \, \exp\!\Big(\beta_{th} \big(T_j - T_{ref}\big)\Big)}_{\text{Potência Estática / Fuga Térmica } (P_{leak})} \Big]
    \label{eq:potencia_ia}
\end{equation}
com lei de escalonamento tensão-frequência (DVFS) e carga térmica total injetada na planta térmica:
\begin{equation}
    V_{core}(f_{clk}) = V_{min} + \kappa_v \left( \frac{f_{clk} - f_{min}}{f_{max} - f_{min}} \right) \quad \implies \quad P_{dyn} \propto f_{clk}^{\,3}
    \label{eq:dvfs_lei}
\end{equation}
\begin{equation}
    q_{gen}(t) = P_{TI}\big(f_{clk}(t), V_{core}(t), T_j(t)\big) + Q_{perdas}(t) + Q_{OISL}
    \label{eq:qgen_total}
\end{equation}

\vspace{2pt}
\noindent\textbf{Parâmetros Nominais:} Fator de atividade computacional $\alpha(t) \in [0{,}1,\; 1{,}0]$; tensão de núcleo $V_{core} \in [0{,}65\text{ V},\; 0{,}85\text{ V}]$; frequência normalizada $f_{clk} \in [0{,}4\,f_{max},\; 1{,}0\,f_{max}]$; sensibilidade térmica de fuga $\beta_{th} \approx 0{,}035\text{ K}^{-1}$; teto rígido de controle $T_j \le 80^\circ\text{C}$.\\
\textbf{Entradas $\to$ Saídas:} Comando DVFS $f_{clk}(t)$, carga $\alpha(t)$ e $T_j(t)$ $\;\longrightarrow\;$ \textbf{$P_{TI}(t)$} e \textbf{Calor total $q_{gen}(t)$ (entrada do Cap.~4)}.
\end{tcolorbox}

\vspace{6pt}

% ==============================================================================
% SÍNTESE DE VARIÁVEIS E ACOPLAMENTO
% ==============================================================================
\section{Síntese de Variáveis, Unidades e Acoplamento entre Capítulos}

\begin{table}[h!]
\centering
\small
\renewcommand{\arraystretch}{1.15}
\begin{tabularx}{\textwidth}{clccX}
    \toprule
    \textbf{Modelo} & \textbf{Variável Principal} & \textbf{Símbolo} & \textbf{Unidade} & \textbf{Papel na Malha de Controle (Cap.~3 $\to$ Cap.~4)} \\
    \midrule
    \textbf{1} & Potência Solar Gerada & $P_{pv}(t)$ & W / kW & Perturbação/entrada energética governada pela órbita. \\
    \textbf{2} & Perdas de Conversão GaN & $Q_{perdas}(t)$ & W / kW & Fonte secundária de calor acoplada à carcaça. \\
    \textbf{3} & Estado de Carga da Bateria & $SoC(t)$ & \% (adim.) & Restrição energética rígida ($SoC \ge 70\%$) no TEMS. \\
    \textbf{4} & Frequência / Tensão DVFS & $f_{clk}, V_{core}$ & GHz / V & Variáveis manipuladas primárias pelo controlador TEMS. \\
    \textbf{4} & Calor Total Injetado & $q_{gen}(t)$ & W / kW & Termo fonte da Equação do Calor e da PINN no Cap.~4. \\
    \bottomrule
\end{tabularx}
\end{table}

\vspace{4pt}

\begin{tcolorbox}[colback=lightgray, colframe=darkgray, title=\color{darkgray}\textbf{FRONTEIRA DE ESCOPO (O QUE FOI DELIBERADAMENTE SIMPLIFICADO NO CAPÍTULO 3)}]
\small
Para preservar o foco em \textbf{Engenharia de Controle e Automação} e manter coerência temporal com a dinâmica térmica orbital (escala de segundos a minutos), os seguintes fenômenos de microescala elétrica foram excluídos da modelagem do Capítulo~3:
\begin{itemize}[leftmargin=1.4em, itemsep=1pt, parsep=0pt]
    \item Chaveamento de alta frequência ($150\text{ kHz}$) e perdas magnéticas em ferritas (equação de Steinmetz);
    \item Análise de estabilidade de pequenos sinais por impedância (critério de Middlebrook) e desligamento por arco em chaves SSPC na escala de microssegundos;
    \item Modelagem eletroquímica de 2ª ordem (\textit{Dual Polarization}) com estimação matricial AUKF Sage--Husa.
\end{itemize}
\end{tcolorbox}

\end{document}
